\chapter{Oscillations}

\mathTopics{\pgRef{chap:vectors}, \pgRef{chap:ellipses}}

\section{Simple Harmonic Motion}

\textbf{Periodic motion} is any specific motion that repeats periodically.

\textbf{Periodic oscillation} is a specific type of periodic motion
where an object moves back and forth in a fix location.

\textbf{Simple harmonic motion} is a type of periodic oscillation
where the net force is only a restoring force
that is a function of the object's distance from equilibrium.

Oscillations can be defined by three parameters.
Let
\begin{itemize}
    \item \(A\) be the \textbf{amplitude},
        which is half the difference between the maximum and minimum values.
        Its SI units depends on what we are measuring,
        e.g., \unit{\metre} for distance, \unit{\radian} for angle,
        \unit{\metre \per \second} for velocity,
        \unit{\ampere} for current, etc.
    \item \(\omega\) be the \textbf{angular frequency},
        which is a measure of how fast the oscillation cycles
        and has SI units of \unit{\radian \per \second}.
    \item \(\phi\) be the \textbf{phase},
        which is the horizontal translation of the cycle,
        and has SI units of \unit{\radian}.
\end{itemize}
Thus, the general form looks like
\begin{equation}
    f(t) = A \cos(\omega t + \phi).
\end{equation}

Several other measurements are related to angular frequency.

\textbf{Rotational frequency} or simply \textbf{frequency} (\(f\))
is the number of oscillations per unit of time,
with SI units of hertz (\unit{\hertz}) or \unit{\per \second}.
The relationship between angular frequency and rotational frequency is
\begin{equation}
    f = \frac{\omega}{2 \pi},
\end{equation}
where \(\omega\) is the angular frequency.

\textbf{Period} (\(T\)) is the time for one complete oscillation,
with SI units of \unit{\second}.
The period \(T\) can be expressed in terms of angular frequency \(\omega\) as
\begin{equation}
    T = \frac{2 \pi}{\omega}.
\end{equation}

\section{Spring Oscillators}

\textbf{Spring oscillators} are composed of a mass
attached to the end of a spring.

Let
\begin{itemize}
    \item \(x\) be the position of the mass.
    \item \(a\) be the acceleration of the mass.
    \item \(m\) be the mass attached to the spring.
    \item \(k\) be the spring constant of the spring.
    \item \(t\) be the time.
\end{itemize}

The position function can be derived from
\begin{equation}
    \begin{aligned}
        F &= m a \\
        -k x &= m \odv[2]{x}{t} \\
        \odv[2]{x}{t} &= - \frac{k}{m} x.
    \end{aligned}
\end{equation}
\begin{subequations}
    We know that the answer when differentiated twice
    is is proportional to the negative of itself,
    so we can guess the solution
    \begin{align}
        x(t) &= A \cos(\omega t + \phi) \\
        v(t) = \odv{x}{t} &= - A \omega \sin(\omega t + \phi) \\
        a(t) = \odv{v}{t} = \odv[2]{x}{t} &= - A \omega^2 \cos(\omega t + \phi),
    \end{align}
    where
    \begin{itemize}
        \item \(A\) is the amplitude,
            which is determined by our initial conditions.
        \item \(\omega\) is the angular frequency as given by the equation
            \begin{equation}
                \omega = \sqrt{\frac{k}{m}}.
            \end{equation}
        \item \(\phi\) is the phase,
            which is determined by our initial conditions.
    \end{itemize}
\end{subequations}

\section{Pendulums}

A \textbf{pendulum} is a swinging rod,
possibly with a mass attached to the end.
A grandfather clock is a good example of a pendulum.

The oscillation of pendulums is determined by gravity.

\subsection{Ideal Pendulums}

For ideal pendulums, i.e., ones with negligible moment of inertia,
the angle function can be derived as following:
\begin{equation} \label{eq:pendulum-forces}
    \begin{aligned}
        F = m a \\
        - m g \sin\theta = m a \\
        - g \sin\theta = a \\
    \end{aligned}
\end{equation}
The arc length \(s\) is given by
\begin{equation}
    \begin{aligned}
        s = \ell \theta. \\
        \therefore a = \ell \odv[2]{\theta}{t}
    \end{aligned}
\end{equation}
Plug back into \cref{eq:pendulum-forces}.
\begin{equation}
    \begin{aligned}
        - g \sin\theta &= \ell \odv[2]{\theta}{t} \\
        \odv[2]{\theta}{t} + \frac{g}{\ell} \sin\theta &= 0
    \end{aligned}
\end{equation}
The \textbf{small-angle approximation} tells us that for sufficiently
small angles
(\qty{0.24}{\radian} gives less than 1\% error),
\begin{gather}
    \sin\theta \approx \theta \\
    \therefore \odv[2]{\theta}{t} + \frac{g}{\ell} \theta \approx 0.
\end{gather}
\begin{subequations}
    If we guess like we did for the spring oscillators, we find that
    \begin{equation}
        \theta(t) = A \cos(\omega t + \phi).
    \end{equation}
    where
    \begin{equation}
        \omega = \sqrt{\frac{g}{\ell}}
    \end{equation}
\end{subequations}

\subsection{Physical Pendulums}

\textbf{Physical pendulums} have non-negligible moment of inertia.

Let \(r\) be the distance between the axis of rotation and center of mass.
The angle function for a physical pendulum is derived as follows:
\begin{equation}
    \begin{aligned}
        \tau &= I \alpha \\
        &= I \odv[2]{\theta}{t} \\
        - r m g \sin\theta &= I \odv[2]{\theta}{t} \\
        \odv[2]{\theta}{t} + \frac{r m g \sin\theta}{I} &= 0 \\
        \odv[2]{\theta}{t} + \frac{r m g \theta}{I} &\approx 0 \\
    \end{aligned}
\end{equation}
\begin{subequations}
    We get
    \begin{equation}
        \theta(t) = A \cos(\omega t + \phi),
    \end{equation}
    where
    \begin{equation}
        \omega = \sqrt{\frac{r m g}{I}}
    \end{equation}
\end{subequations}

\section{Torsion Pendulums}

\textbf{Torsion} is a torque caused by twisting an object around its
axis of rotation.
In a \textbf{torsion pendulum}, an object is hung on a string,
and the string is twisted,
which creates a restorative force of torsion
that causes the object to rotate in an oscillatory manner.
For a torsion pendulum,
the SI units for the spring constant \(k\) is \unit{\newton \per \radian}.

Skipping the differential equations, the derivation for a torsion pendulum is
\begin{subequations}
    \begin{equation}
        \begin{aligned}
            - \tau &= I \alpha \\
            - k \theta &= I \odv[2]{\theta}{t} \\
            - \frac{k}{I} \theta &= \odv[2]{\theta}{t} \\
            \theta(t) &= \theta_\mathrm{max} \cos(\omega t + \phi),
        \end{aligned}
    \end{equation}
    where
    \begin{itemize}
        \item \(\theta_\mathrm{max}\) is the amplitude.
            If the initial angular velocity is 0,
            then \(\theta_\mathrm{max}\) will be the initial angle.
            Elsewise, you will have to solve for \(\theta_\mathrm{max}\)
            using the initial angle and angular velocity.
        \item
            \begin{equation}
                \omega = \sqrt{\frac{k}{I}}
            \end{equation}
            is the angular frequency.
        \item \(\phi\) is the phase shift.
            If the initial angular velocity is 0, \(\phi = 0\).
            Elsewise, you will have to solve for \(\phi\)
            using the initial angle and angular velocity.
    \end{itemize}
\end{subequations}

\section{Resonance}

When something is oscillating,
and an external force is applied at the same frequency as the oscillation,
the amplitude of the oscillation will increase.
This phenomena is known as \textbf{resonance}.
